%% ============================================================================
%% Fichier  : annexes/Ann_A_glossaire.tex
%% Titre    : Glossaire
%% Version  : 1.0 -- novembre 2026
%% ============================================================================
%% Définitions abrégées des notions du cours ; le chapitre de référence est
%% indiqué. Alignée sur l'annexe A des autres cours de l'auteur.
%% ============================================================================

\chapter{Glossaire}
\label{ann:glossaire}

Ce glossaire reprend, sous une forme abrégée, les notions du cours. Les
notions de base (acquisition, chaîne de custody, empreinte, bloqueur
d'écriture) sont définies dans l'UE SEC4031. Le renvoi indique le chapitre
où la notion est présentée~; c'est là qu'il faut lire la définition complète.

\begin{description}
    \item[\$MFT, \$UsnJrnl] Table maîtresse des fichiers et journal des
          modifications du système de fichiers NTFS~; sources d'horodatages
          (Ch.~\ref{chap:for-sys}).
    \item[Anti-forensique] Ensemble des techniques qui visent à empêcher,
          compliquer ou tromper l'investigation (Ch.~\ref{chap:anti-for}).
    \item[Artefact d'exécution] Trace qui prouve qu'un programme a été lancé~:
          Prefetch, Amcache, ShimCache, UserAssist (Ch.~\ref{chap:for-sys}).
    \item[Attribution] Rattachement d'une action à une personne ou à un
          groupe~; elle exige plus qu'une adresse IP (Ch.~\ref{chap:for-osint},
          \ref{chap:rapport}).
    \item[Beaconing] Connexions régulières d'un programme malveillant vers son
          serveur de commande (Ch.~\ref{chap:for-res}).
    \item[CloudTrail] Service de journalisation des appels d'interface d'un
          fournisseur de nuage (Ch.~\ref{chap:for-cloud}).
    \item[Conteneur] Processus isolé qui partage le noyau de l'hôte~; il peut
          disparaître avec ses traces (Ch.~\ref{chap:containers}).
    \item[Constatation et analyse] Le fait vérifiable, avec sa source~;
          l'interprétation, avec un niveau de confiance (Ch.~\ref{chap:rapport}).
    \item[Cryptanalyse] Étude des moyens de mettre en défaut un mécanisme
          cryptographique (Ch.~\ref{chap:cryptanal}).
    \item[Deepfake] Contenu audio ou vidéo synthétique qui imite une personne
          réelle (Ch.~\ref{chap:ia-for}).
    \item[Dossier probatoire] Les pièces du rapport, avec leurs empreintes,
          qui permettent à un expert adverse de tout vérifier
          (Ch.~\ref{chap:rapport}).
    \item[Éphéméralité] Disparition rapide des données dans le nuage ou les
          conteneurs, qui impose de conserver d'abord (Ch.~\ref{chap:for-cloud},
          \ref{chap:containers}).
    \item[Exfiltration] Sortie non autorisée de données d'un système
          (Ch.~\ref{chap:for-res}).
    \item[Firmware (micrologiciel)] Logiciel embarqué d'un équipement, à
          extraire et à analyser dans l'investigation des objets connectés
          (Ch.~\ref{chap:for-iot}).
    \item[Fonction de hachage] Fonction qui associe à des données une empreinte
          de taille fixe, résistante aux collisions (Ch.~\ref{chap:prim-crypt}).
    \item[Hypothèse alternative] Explication concurrente que le rapport doit
          examiner et, s'il le peut, réfuter (Ch.~\ref{chap:rapport}).
    \item[IoT] Objets connectés~: routeurs, caméras, enregistreurs, capteurs
          (Ch.~\ref{chap:for-iot}).
    \item[Locard inversé] Toute tentative d'effacer une trace en produit de
          nouvelles (Ch.~\ref{chap:anti-for}).
    \item[Méthode formelle] Analyse mathématique d'un protocole à l'aide d'un
          modèle et d'un outil de vérification (Ch.~\ref{chap:form-meth}).
    \item[Mobile Money] Service de paiement par téléphone mobile, source
          fréquente de preuves et de fraudes (Ch.~\ref{chap:for-mob}).
    \item[Modèle de responsabilité partagée] Répartition des contrôles et des
          journaux entre le fournisseur de nuage et son client
          (Ch.~\ref{chap:for-cloud}).
    \item[Non-répudiation] Propriété qui empêche l'auteur d'une action de
          la nier de façon crédible (Ch.~\ref{chap:ZKNR}).
    \item[NVR] Enregistreur vidéo en réseau des caméras de surveillance
          (Ch.~\ref{chap:for-iot}).
    \item[OSINT] Collecte et analyse d'informations issues de sources
          légalement accessibles au public (Ch.~\ref{chap:for-osint}).
    \item[Post-quantique] Se dit d'un algorithme conçu pour résister à un
          ordinateur quantique (Ch.~\ref{chap:quantique}).
    \item[Prefetch] Fichiers de Windows qui gardent la trace des programmes
          lancés et de leur nombre d'exécutions (Ch.~\ref{chap:for-sys}).
    \item[Preuve à divulgation nulle de connaissance] Preuve qui convainc
          qu'une affirmation est vraie sans rien révéler d'autre
          (Ch.~\ref{chap:prim-crypt}, \ref{chap:ZKNR}).
    \item[Registre Windows] Base de configuration de Windows, riche en traces
          d'activité (Ch.~\ref{chap:for-sys}).
    \item[Rapport forensique] Livrable en trois niveaux~: synthèse,
          rapport technique, dossier probatoire (Ch.~\ref{chap:rapport}).
    \item[Signature BLS] Schéma de signature qui permet l'agrégation de
          signatures (Ch.~\ref{chap:cas-zknr}).
    \item[SQLite] Moteur de base de données embarqué dans la plupart des
          applications mobiles (Ch.~\ref{chap:for-mob}).
    \item[Timeline] Chronologie unique de toutes les sources, horloges
          corrigées (Ch.~\ref{chap:for-sys}).
    \item[Timestomping] Modification des horodatages d'un fichier pour
          tromper la chronologie (Ch.~\ref{chap:anti-for}).
    \item[Triage assisté] Classement d'un volume de données par pertinence
          probable, sans écarter de pièce sans examen (Ch.~\ref{chap:ia-for}).
    \item[Trilemme CRO, cadre CLO] Tension formalisée entre confidentialité,
          fiabilité et opposabilité, et sa mesure (Ch.~\ref{chap:CRO}~;
          \cite{minka2025cro}).
    \item[Volatility] Cadriciel d'analyse de la mémoire vive
          (Ch.~\ref{chap:for-sys}).
    \item[ZK-NR] Architecture cryptographique en couches pour une
          non-répudiation explicable, publiée en prépublication
          (Ch.~\ref{chap:ZKNR}~; \cite{minka2025zknr}).
\end{description}
